\section{Space Complexity --- Savitch ARM --- Reach}
%%% generated by the blueprint pipeline from TCSlib.Complexity.SpaceComplexity.SavitchARM.Reach
\subsection*{Space Complexity --- Savitch ARM --- Reach: Introduction}

Declarations of the Lean module \texttt{TCSlib.Complexity.SpaceComplexity.SavitchARM.Reach}.

\subsection*{Space Complexity --- Savitch ARM --- Reach: Results}

\begin{definition}[Bounded recursive reachability]
\lean{Complexity.LogProg.SavitchARM.Reach}
\statementsource{Sav70}{
  pdf-c8033d2fd705-p004-b001,
  pdf-c8033d2fd705-p007-b002
}
\label{Complexity.LogProg.SavitchARM.Reach}
\leanok
For a natural number $K$ and a binary relation $R$ on the natural numbers, define
reachability at level $0$ by $u=v$ or $R(u,v)$. At level $i+1$, declare $u$ and $v$
related when there is some $m<K$ such that $u$ is related to $m$ at level $i$ and $m$ is
related to $v$ at level $i$.
\end{definition}

\begin{theorem}[Reachability Along a Chain]
\lean{Complexity.LogProg.SavitchARM.reach_of_chain}
\statementsource{Sav70}{
  pdf-c8033d2fd705-p007-b002,
  pdf-c8033d2fd705-p008-b001
}
\label{Complexity.LogProg.SavitchARM.reach_of_chain}
\leanok
\uses{Complexity.LogProg.SavitchARM.Reach}
\difficulty{5}
Let $K,T$ be natural numbers, let $R$ be a relation on the natural numbers, and let $c$
be a sequence of natural numbers. Define reachability at level $0$ to mean equality or
one $R$-step; at each higher level, two points are reachable if there is an intermediate
number less than $K$ such that each pair of successive points is reachable at the
preceding level. Suppose $R(c_j,c_{j+1})$ for every $j<T$, and $c_j<K$ for every $j\le
T$. Then for all natural numbers $i,j,t$ with $j+t\le T$ and $t\le 2^i$, the points
$c_j$ and $c_{j+t}$ are reachable at level $i$.
\end{theorem}

\begin{theorem}[Reachability Gives Reflexive Transitive Reachability]
\lean{Complexity.LogProg.SavitchARM.reflTransGen_of_reach}
\label{Complexity.LogProg.SavitchARM.reflTransGen_of_reach}
\leanok
\uses{Complexity.LogProg.SavitchARM.Reach}
\difficulty{4}
Let $K$ be a natural number and $R$ a relation on the natural numbers. If $u$ reaches
$v$ at any level under the recursively defined reachability relation—where level $0$
allows equality or one $R$-step and each higher level composes two reachability steps
through an intermediate point below $K$—then $v$ is reachable from $u$ by a finite
sequence of $R$-steps, allowing a sequence of length zero.
\end{theorem}

\begin{theorem}[Configuration Count Bound]
\lean{Complexity.LogProg.SavitchARM.configBound_le}
\label{Complexity.LogProg.SavitchARM.configBound_le}
\leanok
\uses{Turing.FinNDTM, Turing.FinNDTM.configBound}
\difficulty{6}
Let $N$ be a finite nondeterministic Turing machine over the Boolean alphabet, with $k$
work tapes and state set $Q$. For natural numbers $n,E$, if $|Q|+2+9k\le E$, then
\[3(|Q|+1)(n+2)\,3^{k(2(n+2)^E+1)}(2(n+2)^E+1)^k\le 2^{(n+2)^{2E}}.\]
\end{theorem}

\begin{theorem}[Code Width Bound]
\lean{Complexity.LogProg.SavitchARM.codeWidth_le}
\label{Complexity.LogProg.SavitchARM.codeWidth_le}
\leanok
\uses{Complexity.Codec.W, Complexity.Codec.codeK, Complexity.logSpace, Complexity.logSpace_le_succ, Turing.FinNDTM}
\difficulty{4}
Let $N$ be a finite nondeterministic Turing machine over the Boolean alphabet, with $k$
work tapes and state set $Q$. Put $H=3(|Q|+1)$ and $C=2(2(H+3k+1))+2$. For natural
numbers $n,E$, if $2C\le E$, then $C\big((n+2)^E+\log_2(n)+1\big)\le(n+2)^{2E}$, where
$\log_2(n)$ denotes the natural-number logarithm appearing in the definition of
logarithmic space.
\end{theorem}

\begin{theorem}[Polynomial Radius Bound]
\lean{Complexity.LogProg.SavitchARM.radius_le}
\label{Complexity.LogProg.SavitchARM.radius_le}
\leanok
\difficulty{4}
For natural numbers $n,k,c,E$, if $k+c+1\le E$, then $c(n^k+1)\le(n+2)^E$.
\end{theorem}

\begin{theorem}[Language membership and bounded configuration reachability]
\lean{Complexity.LogProg.SavitchARM.mem_iff_reach}
\label{Complexity.LogProg.SavitchARM.mem_iff_reach}
\leanok
\uses{Complexity.Codec.StepRel, Complexity.Codec.acceptLang, Complexity.Codec.codeK, Complexity.Codec.initLang, Complexity.Codec.mem_iff_codeChain, Complexity.Codec.mem_of_codePath, Complexity.LogProg.SavitchARM.Reach, Complexity.LogProg.SavitchARM.reach_of_chain, Complexity.LogProg.SavitchARM.reflTransGen_of_reach, Complexity.logSpace, Turing.FinNDTM, Turing.FinNDTM.DecidesInSpace, Turing.FinNDTM.configBound, Turing.pairEncode}
\difficulty{5}
Let $N$ be a nondeterministic Turing machine with $k$ work tapes and a finite state set,
and suppose $N$ decides a language $L$ within space bound $\mathrm{sp}$. Let $x$ be a
Boolean string, and let $s,w\in\mathbb{N}$. Write $q$ for the number of states of $N$,
set $d=2\bigl(3(q+1)+3k+1\bigr)+2$, and define $B=3(q+1)(|x|+2)3^{k(2s+1)}(2s+1)^k$. If
$\mathrm{sp}(|x|)\le s$, $d\bigl(s+\lfloor\log_2|x|\rfloor+2\bigr)\le w$, and $B\le
2^w$, then $x\in L$ if and only if there are natural numbers $u,v<2^w$ whose encodings
with $x$ and $s$ represent, respectively, the initial configuration and an accepting
configuration of $N$, and $v$ is reachable from $u$ through the encoded one-step
configuration relation at recursive depth $w$, with intermediate codes less than $2^w$.
\end{theorem}
